% TOP_PROBLEM: {"id": 2307079, "title": "Research Problems in Function Theory — Problem 7.79", "category_id": 8, "category": {"id": 8, "name": "analysis", "display_name": "Analysis", "description": "Limits, continuity, calculus, and function theory.", "slug": "analysis", "order_index": 8, "created_at": "2026-07-31T15:26:25.671Z"}, "status": "open"}
% FIRST_POSTED: null
% ENTRY_KIND: statement_only
% Imported from: batch06.tex; UnsolvedMath ID 2307079
% Compile twice: lualatex 2307079.tex
\documentclass[11pt]{article}
\usepackage{fontspec}
\setmainfont{Latin Modern Roman}
\usepackage{amsmath,amssymb,mathtools,unicode-math}
\setmathfont{Latin Modern Math}
\usepackage{xurl,hyperref}
\hypersetup{hidelinks,pdftitle={UnsolvedMath Problem 2307079}}
\newcommand{\sref}[2]{\hyperlink{source:#1:#2}{[S#2]}}
\newcommand{\cataloguescope}[1]{\paragraph{Mathematical question.}#1}
\begin{document}
% UnsolvedMath ID 2307079; source catalogue code AMR-022-7079
\setcounter{section}{2307078}
\section{Research Problems in Function Theory — Problem 7.79}\label{2307079}
{\small\sffamily UnsolvedMath ID 2307079 · Analysis}
\subsection{Definitions and mathematical statement}

\paragraph{Shared notation.}$\mathbb N=\{1,2,\ldots\}$, and $\mathbb Z,\mathbb Q,\mathbb R,\mathbb C$ denote the integers, rationals, reals and complex numbers. Any other conventions are specified locally.


Let $G\subset\mathbb C$ be a domain and $f$ holomorphic there. Define $M(f)$ to consist of points with a neighborhood $N\subset G$ on which $\{f^{(n)}:n\ge1\}$ is normal. Normality means that every sequence has a subsequence converging locally uniformly in the spherical metric, with holomorphic/meromorphic limit or identically infinity as appropriate. Define $M_!(f)$ in the same way using $\{f^{(n)}/n!:n\ge1\}$. For meromorphic $f$ use the same spherical convention. A set is simply connected here if it is connected and every loop contracts within it.
\paragraph{Statement recorded in the supplied source.}
Characterize the possible sets $M(f)$ beyond their being open. Must $M(f)$ be connected? If $G$ is simply connected, must $M(f)$ be simply connected? Can it be an arbitrary prescribed open subset of $G$? Answer the analogous questions for meromorphic $f$, and for the factorial-normalized derivative family defining $M_!(f)$. \sref{2307079}{2}
\subsection{Short English statement}
Determine which open sets can be the region where a function’s derivatives form a normal family, including connectivity and the meromorphic and factorial-normalized variants.
\subsection{Sources}
\begin{description}
\item[\hypertarget{source:2307079:1}{S1}] ULAM.AI and the UnsolvedMath Contributors, UnsolvedMath: A Curated Collection of Open Mathematics Problems, record 2307079 (AMR-022-7079). CC BY 4.0. \url{https://huggingface.co/datasets/ulamai/UnsolvedMath}.
\item[\hypertarget{source:2307079:2}{S2}] Walter K. Hayman and Edward F. Lingham, Research Problems in Function Theory (2018), Problem 7.79, its complete update and relevant preceding definitions. \url{https://arxiv.org/abs/1809.07200}.
\end{description}
\label{end:2307079}
\end{document}
